\section{局限、风险与合规}
\label{sec:limits}

\paragraph{评测规模。}真实年报抽取只覆盖三家公司，研报质检只覆盖五家公司、七份交付稿，每份只评一次。自造样例上的满分指标只能说明流程正确。我们正在收集约 15 家、覆盖消费、制造、周期、科技、公用事业等行业与沪深北三地版式的年报与季报，以及针对其中 8～10 家的卖方研报人工埋错集（计划 150～200 处，其中约 20\% 为“正确但容易误报”的断言），用于更可靠地评估抽取与纠错。

\paragraph{非对照比较。}研报得分的演进同时改变了多个因素，不是消融实验。我们计划在扩充评测集后，分别去掉资料搜集先行、复核会话与数字锁、深度专题三项设计，测量各自的贡献。

\paragraph{评审模型的局限。}评审与生成模型属于同一模型家族，可能偏好同族生成的文本。重复评审的总分极差为 7～8 分，接近交付线附近的得分差。整节删除类缺陷的判出率偏低，虽然主要原因是劣化不彻底，但也暴露了“是否存在”类条目容易被零散内容满足的问题。人工复核只覆盖评审判“不通过”的条目，无法估计漏判率。

\paragraph{闸门是数量下限。}字数、页数、图表数、来源数等闸门只能阻止模型偷懒，不能认证质量；一份满足全部数量闸门的研报仍可能论证薄弱，这部分只能靠盲评与人工复核兜底。

\paragraph{护栏的边界。}护栏核对的是“有出处”与“口径明确的指标绑定”，不核对预测是否合理、外部来源是否权威、分析判断是否正确。预测数字之所以能通过护栏，是因为它们出现在计算产物中，而这些产物本身由模型起草的假设驱动；假设的合理性要靠假设表中的依据、反向 DCF 的对照以及人工复核清单来约束。

\paragraph{合规与投资建议。}所有估值参数均为起草的假设，参考股价为外部数据，估值结论高度依赖这些假设。系统生成的评级由固定规则从估值结果得出，研报统一附“不构成投资建议”的免责声明，不提供加仓、减仓等操作建议。卖方研报与年报原件的版权归原作者所有：评测集中的卖方研报只用于内部评测，仓库只保存我们抽出的断言与标签。第三方代码的参照关系、版本与许可证在 \file{THIRD\_PARTY.md} 中逐项列明。
